% koide_ckm_prl.tex -- PRL letter fusing Casimir-Koide + CKM/PMNS results
\documentclass[aps,prl,twocolumn,showpacs,preprintnumbers,superscriptaddress]{revtex4-2}

\usepackage{amsmath,amssymb,amsthm}
\usepackage{graphicx}
\usepackage{hyperref}
\usepackage{booktabs}

\newcommand{\SU}{\mathrm{SU}}
\newcommand{\UU}{\mathrm{U}}
\newcommand{\Tr}{\mathrm{Tr}}
\newcommand{\adj}{\mathrm{adj}}
\newcommand{\diag}{\mathrm{diag}}
\newcommand{\MeV}{\,\mathrm{MeV}}
\newcommand{\GeV}{\,\mathrm{GeV}}
\newcommand{\TeV}{\,\mathrm{TeV}}
\newcommand{\Nc}{N_{\!c}}
\newcommand{\Nf}{N_{\!f}}

\newtheorem{theorem}{Theorem}

\begin{document}

\title{The Koide ratio as a Casimir invariant: generation uniqueness
and a $1/N$ flavour expansion from $\SU(3)_F$}

\author{A.~Rivero}
\email{rivero@unizar.es}
\affiliation{Departamento de F\'{\i}sica Te\'orica, Universidad de
Zaragoza, 50009 Zaragoza, Spain}

\date{\today}

\begin{abstract}
The charged-lepton mass ratio
$K \equiv (m_e + m_\mu + m_\tau)/(\sqrt{m_e}+\sqrt{m_\mu}+\sqrt{m_\tau})^2
= 0.666661$ agrees with $2/3$ to $0.001\%$.
We prove that $K = 2/3$ is the mean-square weight
$\langle J_3^2\rangle$ of the principal $\mathfrak{sl}(2)$
subalgebra of $\mathfrak{sl}(N)$ in the fundamental representation,
and that $K_N = \langle J_3^2\rangle$ holds if and only if
$N(N^2-1) = 24$, whose unique positive integer solution is $N = 3$.
This Casimir--Koide identity is an algebraic theorem.
Separately, we observe that the same Cartan angle
$\delta = 2/9$ parametrises all six CKM and PMNS mixing
parameters as rational functions of $N$, matching data
with $\chi^2/\text{dof} = 6.2/5$ ($p = 0.28$) when
$|V_{us}|$ fixes~$\delta$.
These mixing relations are empirical patterns admitting a
compact $1/N$ representation; they are not derived from
confining dynamics.
\end{abstract}

\maketitle


%======================================================================
% INTRODUCTION
%======================================================================

The charged-lepton mass ratio
\begin{equation}
  K \;\equiv\; \frac{m_e + m_\mu + m_\tau}
  {\bigl(\sqrt{m_e}+\sqrt{m_\mu}+\sqrt{m_\tau}\bigr)^2}
  \;=\; 0.666661
  \label{eq:Kdef}
\end{equation}
was noted by Koide~\cite{Koide:1983} to agree with $2/3$ to
one part in $10^5$.
In the parametrisation
$\sqrt{m_n} = \mu\bigl(1 + \sqrt{2}\cos(\delta + 2\pi n/3)\bigr)$
of Brannen~\cite{Brannen:2006}, Foot's geometric
condition~\cite{Foot:1994} $\theta = \pi/4$ is equivalent to
$K = 2/3$, and the mass hierarchy is encoded in a single
Cartan angle~$\delta$.
Numerically $\delta = 0.22222 \pm 0.00001$~rad,
consistent with $\delta = 2/9$.
Despite numerous models~\cite{Koide:1983,Sumino:2009,Rivero:2024},
no derivation of the value $2/3$ from first principles exists
without introducing it as an input.

We report two results.
First, $K = 2/3$ is a representation-theoretic invariant:
the mean-square weight $\langle J_3^2\rangle$ of the principal
$\mathfrak{sl}(2) \subset \mathfrak{sl}(N)$~\cite{Kostant:1959}
in the fundamental representation, with the equality
$K_N = \langle J_3^2\rangle$ holding if and only if $N = 3$.
This is an algebraic theorem.
Second, we observe that the same angle $\delta = 2/9$
parametrises all six CKM and PMNS mixing parameters as rational
functions of $N$, with no free parameters beyond $N = 3$.

%======================================================================
% CASIMIR-KOIDE IDENTITY
%======================================================================

\textit{The Casimir--Koide identity.}---Let
$\hat{z}_i = z_0 + z_i$ ($i = 1,\ldots,N$) be the signed
charges of a $\UU(N)$ family symmetry, with $z_0$ the $\UU(1)$
generator and $z_i$ the $\SU(N)$ Cartan components.
Canonical normalisation gives
\begin{equation}
  \sum_{i=1}^{N} z_i = 0\,,\qquad
  \sum_{i=1}^{N} z_i^2 = \frac{1}{2}\,,\qquad
  z_0 = \frac{1}{\sqrt{2N}}\,.
  \label{eq:u3_norm}
\end{equation}

\begin{theorem}[Generalised Koide ratio]
\label{thm:KN}
For the mass law $m_i = \mu\,\hat{z}_i^2$ and the ratio
$K_N \equiv \sum\hat{z}_i^2 / (\sum\hat{z}_i)^2$,
one has $K_N = 2/N$ independently of the Cartan direction.
\end{theorem}

\noindent\textit{Proof.}---Using Eq.~\eqref{eq:u3_norm}:
$\sum\hat{z}_i^2 = Nz_0^2 + \sum z_i^2 = 1/2 + 1/2 = 1$,
$\sum\hat{z}_i = Nz_0 = \sqrt{N/2}$,
so $K_N = 1/(N/2) = 2/N$.~$\square$

For $N = 3$: $K_3 = 2/3$, recovering the Koide value as
an algebraic identity of $\UU(3)$.
This is the charge-squared uniqueness
theorem~\cite{Koide:1983,Rivero:2024}: among monomial mass
functions $m \propto \hat{z}^n$, only $n = 2$ gives a
direction-independent ratio.

The principal $\mathfrak{sl}(2)$ of a simple Lie algebra
$\mathfrak{g}$ is the unique (up to conjugacy)
$\mathfrak{sl}(2)$ subalgebra for which the fundamental
representation of $\mathfrak{g}$ is
irreducible~\cite{Kostant:1959}.
For $\mathfrak{g} = \mathfrak{sl}(N)$, the fundamental
$\mathbf{N}$ is the spin-$j$ representation with
$j = (N{-}1)/2$.
The mean-square weight is
\begin{equation}
  \langle J_3^2\rangle
  \;\equiv\; \frac{1}{N}\sum_{m=-j}^{j} m^2
  \;=\; \frac{N^2 - 1}{12}\,.
  \label{eq:J3sq}
\end{equation}
Setting $K_N = \langle J_3^2\rangle$:
\begin{equation}
  \frac{2}{N} = \frac{N^2 - 1}{12}
  \quad\Longleftrightarrow\quad
  \boxed{N(N^2-1) = 24}\,.
  \label{eq:diophantine}
\end{equation}
Factoring: $N(N{-}1)(N{+}1) = 2\times 3\times 4 = 24$.
Since $3\times 4\times 5 = 60 > 24$, the unique positive
integer solution is $N = 3$.
Equivalently, the factored form
$(N-3)(N^2 + 3N + 8) = 0$ has no other real roots
(discriminant $9 - 32 < 0$).

The Diophantine condition~\eqref{eq:diophantine} is the
theorem; the following is a corollary.
At $N = 3$ and only $N = 3$,
the representation dimension $2j+1 = N$ coincides with
$\dim\,\mathfrak{sl}(2) = 3$, so the fundamental of $\SU(3)$
\emph{is} the adjoint of its principal $\mathfrak{sl}(2)$.
This gives the structural coincidence
$C_2/N = C_2/3 = \langle J_3^2\rangle = K = 2/3$.

%======================================================================
% CABIBBO = CARTAN
%======================================================================

\textit{Cabibbo angle = Cartan angle.}---Koide's original
composite model~\cite{Koide:1983} predicted
\begin{equation}
  \tan\theta_C = \frac{\sqrt{3}\,(\sqrt{m_\mu}-\sqrt{m_e})}
  {2\sqrt{m_\tau}-\sqrt{m_\mu}-\sqrt{m_e}}\,.
  \label{eq:cabibbo}
\end{equation}
In the Koide parametrisation, the right-hand side of
Eq.~\eqref{eq:cabibbo} is \emph{algebraically identical}
to $\tan\delta$, for any masses satisfying $K = 2/3$.
Therefore $\theta_C = \delta$.
Numerically $\theta_C = 0.22223$~rad, matching the PDG
value~\cite{PDG:2024} $|V_{us}| = 0.2250 \pm 0.0007$
to~$0.9\%$.

The identification $\delta = 2/N^2$ with $N = 3$
expresses the Cabibbo angle as $|V_{us}| = 2/9$.
Combining with the Gatto--Sartori--Tonin
relation~\cite{GST:1968} $|V_{us}| \approx \sqrt{m_d/m_s}$
gives a cross-sector prediction:
\begin{equation}
  \frac{m_d}{m_s} = \delta^2 = \frac{4}{81} = 0.0494\,.
  \label{eq:md_ms}
\end{equation}
The PDG value $m_d/m_s = 0.050 \pm 0.004$ agrees to $1.2\%$.

%======================================================================
% CKM FROM DEMOCRATIC RESOLVENT
%======================================================================

\textit{CKM from the democratic resolvent.}---The democratic
matrix $J = \mathbf{1}\mathbf{1}^T$ has eigenvalues $N$
(democratic direction) and~$0$ ($(N{-}1)$-fold degenerate
breaking plane).
The resolvent of the zero-eigenvalue subspace at
$z = 1 + \delta$ yields $A = 1/(1+\delta)$---a suggestive
algebraic connection to the Wolfenstein~$A$ parameter, not a
derivation from $\SU(3)_F$ dynamics:
\begin{equation}
  A = \frac{1}{1+\delta} = \frac{N^2}{N^2+2} = \frac{9}{11}\,.
  \label{eq:A}
\end{equation}
The resulting zero-parameter parametrisation gives
\begin{equation}
  |V_{us}| = \frac{2}{9}\,,\qquad
  |V_{cb}| = A\lambda^2 = \frac{4}{99}\,,\qquad
  A = \frac{9}{11}\,.
  \label{eq:ckm}
\end{equation}

\begin{table}[b]
\caption{CKM parametrisation vs.\ PDG~2024~\cite{PDG:2024}.}
\label{tab:ckm}
\begin{ruledtabular}
\begin{tabular}{lccc}
  & $1/N$ value & PDG & Deviation \\
\hline
$|V_{us}|$ & $2/9 = 0.2222$ & $0.2250 \pm 0.0007$ & $1.2\%$ \\
$|V_{cb}|$ & $4/99 = 0.0404$ & $0.0422 \pm 0.0008$ & $4.3\%$\footnotemark[1] \\
$A$ & $9/11 = 0.818$ & $0.826 \pm 0.013$ & $1.0\%$ \\
\end{tabular}
\end{ruledtabular}
\footnotetext[1]{The PDG average uses the inclusive $|V_{cb}|$.
The exclusive determination $0.0395 \pm 0.0008$ gives only
$2.2\%$ deviation.}
\end{table}

%======================================================================
% PMNS FROM 1/N
%======================================================================

\textit{PMNS from the $1/N$ expansion.}---We observe that the PMNS
mixing angles are deviations from maximal mixing ($\pi/4$),
parametrised by~$\delta$:
\begin{align}
  \theta_{12} &= \frac{\pi}{4} - \delta + \frac{\delta^2}{N}
  = \frac{\pi}{4} - \frac{2}{9} + \frac{4}{243}\,,
  \label{eq:t12} \\
  \theta_{23} &= \frac{\pi}{4} + \frac{\delta}{N}
  = \frac{\pi}{4} + \frac{2}{27}\,,
  \label{eq:t23} \\
  \theta_{13} &= \frac{2j(j{+}1)}{N^3}
  = \frac{4}{27}\;\text{rad}\,,
  \label{eq:t13}
\end{align}
where $j = (N{-}1)/2 = 1$.
This extends the quark--lepton complementarity of
Raidal~\cite{Raidal:2004} and the quark Koide phases
of Zenczykowski~\cite{Zenczykowski:2012}.

\begin{table}[b]
\caption{PMNS parametrisation vs.\ NuFIT~5.2.
The correction coefficients $(-1,+1/3,+2/3)$ multiplying
$\delta$ in Eqs.~\eqref{eq:t12}--\eqref{eq:t13} are
traceless and unique among denominators~$\leq 3$.}
\label{tab:pmns}
\begin{ruledtabular}
\begin{tabular}{lccc}
  & $1/N$ value & NuFIT~5.2 & Deviation \\
\hline
$\sin^2\theta_{12}$ & $0.300$ & $0.304 \pm 0.012$ & $1.3\%$ \\
$\sin^2\theta_{23}$ & $0.574$ & $0.573 \pm 0.018$ & $0.1\%$ \\
$\sin^2\theta_{13}$ & $0.0218$ & $0.0222 \pm 0.0007$ & $1.9\%$ \\
\end{tabular}
\end{ruledtabular}
\end{table}

Scanning $N = 2,\ldots,10$ in
Eqs.~\eqref{eq:t12}--\eqref{eq:t13},
only $N = 3$ gives $\chi^2 < 1$:
$\chi^2 = 0.5$ at $N = 3$ versus $604$ ($N = 2$)
and $218$ ($N = 4$).

%======================================================================
% 1/N TABLE
%======================================================================

\textit{The $1/N$ expansion.}---All parameters are rational
functions of $N$ (Table~\ref{tab:1overN}).
The lepton mass hierarchy, the CKM matrix, and the PMNS
matrix all admit a compact $1/N$ representation at $N = 3$.
These relations are empirical observations that admit a
$1/N$ expansion; they are not derived from the confining
dynamics.

\begin{table}[t]
\caption{Flavour parameters as rational functions of~$N$
with $j = (N{-}1)/2$, evaluated at $N = 3$.}
\label{tab:1overN}
\begin{ruledtabular}
\begin{tabular}{llcc}
Parameter & Formula & $N{=}3$ & Order \\
\hline
$\delta$ & $2/N^2$ & $2/9$ & $O(1/N^2)$ \\
$K$ & $j(j{+}1)/N$ & $2/3$ & $O(1/N)$ \\
$|V_{us}|$ & $2/N^2$ & $2/9$ & $O(1/N^2)$ \\
$A$ & $N^2/(N^2{+}2)$ & $9/11$ & $1{-}O(1/N^2)$ \\
$\theta_{12}{-}\pi/4$ & $-2/N^2 + 4/N^5$ & $-50/243$ & $O(1/N^2)$ \\
$\theta_{23}{-}\pi/4$ & $2/N^3$ & $2/27$ & $O(1/N^3)$ \\
$\theta_{13}$ & $2j(j{+}1)/N^3$ & $4/27$ & $O(1/N^3)$ \\
\end{tabular}
\end{ruledtabular}
\end{table}

The CKM elements scale as $\delta^{|i-j|} = O(1/N^{2|i-j|})$:
the quark mixing hierarchy is a power series in $1/N^2$.
The PMNS corrections scale as $O(1/N^{|i-j|-1})$:
PMNS is near-maximal because its corrections are $O(1/N)$,
while CKM is hierarchical because its elements are $O(1/N^2)$.
This structural difference---not fine-tuning---explains
$\theta_{\text{CKM}} \sim 0.22$ versus
$\theta_{\text{PMNS}} \sim O(1)$.

%======================================================================
% DISCUSSION
%======================================================================

\textit{Discussion.}---The statistical assessment depends on
the treatment of~$\delta$.
(a)~With $|V_{us}|$ as input (fixing $\delta$ from data),
the remaining five observables have
$\chi^2/\text{dof} = 6.2/5$ ($p = 0.28$), with pulls:
$|V_{cb}|$ ($-2.3\sigma$), $A$ ($-0.6\sigma$),
$\sin^2\theta_{12}$ ($-0.3\sigma$),
$\sin^2\theta_{23}$ ($+0.04\sigma$),
$\sin^2\theta_{13}$ ($-0.6\sigma$).
(b)~With $\delta = 2/9$ as an exact prediction, all six
observables give $\chi^2/\text{dof} \approx 24/6$,
driven by the $|V_{us}|$ pull of $-4.2\sigma$.
The tension between $2/9 = 0.2222$ and the PDG value
$0.2250$ may be resolved by the instanton-generated
potential (see below), whose minimum at $\delta = 0.2225$
reduces the $|V_{us}|$ pull to $-3.6\sigma$.

The physical origin is a confining $\SU(3)_F$ family gauge
theory~\cite{Rivero:2024} whose Koide formula is the
Casimir identity $K = \langle J_3^2\rangle$, holding iff
$N(N^2{-}1) = 24$ iff $N = 3$.
The operator connecting the two mass branches of $\SU(3)$
SQCD~\cite{Seiberg:1995}---the direct ($\adj(M) \propto m$)
and inverse ($M \propto 1/m$)---is the spin-2 component of
$\mathrm{End}(V)$ under the principal
$\mathfrak{sl}(2)$~\cite{Kostant:1959},
Kostant's second exponent of $\mathfrak{sl}(3)$.

Discrete flavour symmetries ($A_4$, $S_3$) also produce
mixing patterns with few parameters~\cite{Altarelli:2010,King:2015},
but require additional group-theoretic input beyond $N$;
the present parametrisation uses $N$ alone.

The critical experimental test is the solar angle.
Our parametrisation gives $\sin^2\theta_{12} = 0.300$, which is
$1.3\%$ below the NuFIT central value $0.304$ and $1.1\sigma$
below the first JUNO measurement
$\sin^2\theta_{12} = 0.309 \pm 0.009$~\cite{JUNO:2025}.
Continued JUNO data-taking will reach
$\pm 0.003$ precision, providing a decisive test.

A candidate mechanism for the Cartan angle is the minimum of
the instanton-generated potential
$V(\delta) = -\cos(3\delta) + (1/\pi)\cos(6\delta)$,
which gives $\delta = 0.2225$ ($0.12\%$ from $2/9$).
The coefficient $1/\pi$ is conjectural; it would follow from
a one-loop instanton determinant in the confining
$\SU(3)_F$ theory if the functional measure produces a
factor of $1/\pi$ per instanton.
Two quantities remain unpredicted: the CKM and PMNS
CP-violating phases.

\begin{acknowledgments}
The author thanks M.~Porter for identifying Seiberg duality
as the mechanism for the sBootstrap during sustained
discussions (2011--2024),
C.~Brannen for the circulant matrix parametrisation
of the Koide formula, and the late M.~Sheppeard (Kea) for
early contributions connecting Koide to representation theory.
During the preparation of this work the author used Claude
(Anthropic) for algebraic verification, numerical checks,
and drafting assistance.
The author reviewed and edited all content and takes full
responsibility for the publication.
\end{acknowledgments}


\begin{thebibliography}{99}

\bibitem{Koide:1983}
  Y.~Koide,
  Phys.\ Lett.\ B \textbf{120}, 161 (1983).

\bibitem{Rivero:2024}
  A.~Rivero,
  Eur.\ Phys.\ J.\ C \textbf{84}, 1058 (2024),
  arXiv:2407.05397.

\bibitem{Kostant:1959}
  B.~Kostant,
  Amer.\ J.\ Math.\ \textbf{81}, 973 (1959).

\bibitem{Sumino:2009}
  Y.~Sumino,
  JHEP \textbf{05}, 075 (2009),
  arXiv:0812.2103.

\bibitem{Brannen:2006}
  C.A.~Brannen,
  \texttt{hep-ph/0606057} (2006).

\bibitem{Foot:1994}
  R.~Foot,
  arXiv:hep-ph/9402242 (1994).

\bibitem{Seiberg:1995}
  N.~Seiberg,
  Nucl.\ Phys.\ B \textbf{435}, 129 (1995),
  arXiv:hep-ph/9411149.

\bibitem{PDG:2024}
  S.~Navas \textit{et al.} (Particle Data Group),
  Phys.\ Rev.\ D \textbf{110}, 030001 (2024).

\bibitem{Raidal:2004}
  M.~Raidal,
  Phys.\ Rev.\ Lett.\ \textbf{93}, 161801 (2004),
  arXiv:hep-ph/0404046.

\bibitem{Zenczykowski:2012}
  P.~Zenczykowski,
  Phys.\ Rev.\ D \textbf{86}, 117303 (2012),
  arXiv:1210.4125.

\bibitem{GST:1968}
  R.~Gatto, G.~Sartori, and M.~Tonin,
  Phys.\ Lett.\ B \textbf{28}, 128 (1968).

\bibitem{Altarelli:2010}
  G.~Altarelli and F.~Feruglio,
  Rev.\ Mod.\ Phys.\ \textbf{82}, 2701 (2010),
  arXiv:0905.3329.

\bibitem{King:2015}
  S.~F.~King,
  J.\ Phys.\ G \textbf{42}, 123001 (2015),
  arXiv:1510.02091.

\bibitem{JUNO:2025}
  A.~Abusleme \textit{et al.} (JUNO Collaboration),
  arXiv:2511.14593 (2025).

\end{thebibliography}

\end{document}
